\documentclass[aspectratio=169,10pt,spanish]{beamer}

\input{../../../_preambulo_beamer.tex}
\input{../../../_comandos_beamer.tex}

\selectlanguage{spanish}

\title{MA1001B - Modelación Estadística para la Toma de Decisiones}
\subtitle{Lección 45: Comparaciones múltiples Tukey HSD}
\author{}
\institute{Tecnol\'ogico de Monterrey}
\date{}

\begin{document}

\begin{frame}[plain]
  \vspace{1.2cm}
  {\Large\bfseries MA1001B - Modelación Estadística para la Toma de Decisiones\par}
  \vspace{0.7cm}
  {\large Lección 45: Comparaciones múltiples Tukey HSD\par}
  \vspace{0.7cm}
  {\color{TecRojo}\rule{\textwidth}{1.2pt}}
  \vspace{0.8cm}

  Facultad MA1001B\\[0.3cm]
  Periodo\\[0.3cm]
  Tecnol\'ogico de Monterrey
\end{frame}

\begin{frame}{La pregunta por pares después de $F$}
  \begin{alertblock}{Pregunta guía}
    Tras un $F$ significativo de tres grupos, ¿qué pares de método de
    empaque difieren, y cómo puede esa familia de pruebas mantener un error
    familiar de 5 por ciento?
  \end{alertblock}

  \vspace{0.6em}
  Al final de esta lección, usted puede:
  \begin{itemize}
    \item correr las tres pruebas $t$ por pares;
    \item leer $p$-valores ajustados Tukey e IC simultáneos;
    \item nombrar los pares que difieren a $\alpha=0.05$ familiar;
    \item explicar por qué las pruebas sin ajuste inflan el error familiar.
  \end{itemize}

  \vfill
  \scriptsize Todos los tiempos de ciclo usados hoy son sintéticos. No se usan datos reales de empresa.
\end{frame}

\begin{frame}{Tres pares, una familia}
  \small
  \begin{center}
    \begin{tabular}{lrr}
      \toprule
      \textbf{Par} & \textbf{Diferencia de medias} & \textbf{$p$ sin ajuste} \\
      \midrule
      Estándar $-$ Guiado & 1.215833 & 0.287762 \\
      Estándar $-$ Automatizado & 5.716390 & 0.000071 \\
      Guiado $-$ Automatizado & 4.500557 & 0.000233 \\
      \bottomrule
    \end{tabular}
  \end{center}

  \vspace{0.5em}
  \begin{exampleblock}{Error familiar}
    Tres pruebas a $\alpha=0.05$ no son una decisión de 5 por ciento.
    El error familiar es la chance de al menos un par falso en la familia.
  \end{exampleblock}
\end{frame}

\begin{frame}[fragile]{Paso 1: $t$ por pares sin ajuste}
  \begin{columns}[T]
    \begin{column}{0.40\textwidth}
      \small
      \begin{lstlisting}
t, p = ttest_ind(a, b,
                 equal_var=True)
reject = p < 0.05
      \end{lstlisting}

      \begin{block}{Salida verificada}
        Par Guiado $p=0.287762$ retener\\
        Pares Auto $p=0.000071$, $0.000233$\\
        Tres pruebas, sin corrección
      \end{block}
    \end{column}
    \begin{column}{0.60\textwidth}
      \centering
      \includegraphics[width=\textwidth]{../figuras/tukey_hsd_01_pares_sin_ajuste.png}
      \vspace{0.2em}
      \scriptsize Diferencias por pares sin ajuste del Paso 1
    \end{column}
  \end{columns}
\end{frame}

\begin{frame}{Tukey HSD mantiene la familia en 5 por ciento}
  \small
  \begin{center}
    \begin{tabular}{lrrr}
      \toprule
      \textbf{Par} & $\mathbf{p_{\text{adj}}}$ & \textbf{IC Tukey} & \textbf{¿Rechazar?} \\
      \midrule
      Auto vs Guiado & 0.000789 & $(1.79,\ 7.21)$ & Sí \\
      Auto vs Estándar & 0.000033 & $(3.00,\ 8.43)$ & Sí \\
      Guiado vs Estándar & 0.521330 & contiene $0$ & No \\
      \bottomrule
    \end{tabular}
  \end{center}

  \vspace{0.4em}
  Automatizado difiere de ambos métodos. Guiado y Estándar no se
  distinguen. El IC que contiene 0 es la decisión de retener.
\end{frame}

\begin{frame}[fragile]{Paso 2: Intervalos simultáneos Tukey}
  \begin{columns}[T]
    \begin{column}{0.40\textwidth}
      \small
      \begin{lstlisting}
tukey = pairwise_tukeyhsd(
  y, groups, alpha=0.05
)
      \end{lstlisting}

      \begin{block}{Salida verificada}
        Guiado vs Estándar se retiene\\
        $p_{\text{adj}}=0.521330$\\
        El IC contiene 0
      \end{block}
    \end{column}
    \begin{column}{0.60\textwidth}
      \centering
      \includegraphics[width=\textwidth]{../figuras/tukey_hsd_02_tukey_hsd.png}
      \vspace{0.2em}
      \scriptsize Intervalos Tukey HSD del Paso 2
    \end{column}
  \end{columns}
\end{frame}

\begin{frame}{Paso 3: Las pruebas sin ajuste inflan el error familiar}
  \begin{columns}[T]
    \begin{column}{0.42\textwidth}
      \small
      5000 experimentos nulos, todas las medias verdaderas iguales.

      \vspace{0.4em}
      FWER por pares sin ajuste $=0.127800$.

      \vspace{0.4em}
      FWER Tukey HSD $=0.049000$ frente a $\alpha=0.05$.

      \vspace{0.5em}
      \textbf{Límite:} las pruebas $t$ por pares sin corrección inflan el
      error familiar. La Lección 46 chequea el supuesto de varianzas
      iguales que ambos procedimientos aún necesitan.
    \end{column}
    \begin{column}{0.58\textwidth}
      \centering
      \includegraphics[width=\textwidth]{../figuras/tukey_hsd_03_error_familiar.png}
      \vspace{0.2em}
      \scriptsize Error familiar empírico del Paso 3
    \end{column}
  \end{columns}
\end{frame}

\begin{frame}{Verificación de conceptos}
  \begin{enumerate}
    \item ¿Cuántas pruebas por pares existen para $k=3$ métodos?
    \item ¿Por qué el error familiar no es lo mismo que el $\alpha=0.05$ de
      una prueba?
    \item ¿Qué par retiene Tukey HSD en este experimento de empaque?
    \item En 5000 experimentos nulos, ¿por qué $0.127800$ es un problema?
  \end{enumerate}

  \vspace{0.6em}
  \begin{block}{Conexión con la Lección 46}
    La sesión puede continuar directamente con gráficas de residuos de
    ANOVA, la prueba de Levene y la desigualdad de varianzas.
  \end{block}

  \vfill
  \scriptsize Fuente: OpenStax, \textit{Introductory Business Statistics 2e},
  Secciones 12.2--12.3, práctica post-hoc por pares. CC BY-NC-SA.
\end{frame}

\appendix



\end{document}
